\subsection{希尔伯特空间的对偶}

定理 3. --- 设 E 为希尔伯特空间。对每个 \(x \in E\) ，设 \(x^{*}\) 为 E 上的连续线性型 \(y \mapsto \langle x | y \rangle\) ；映射 \(x \mapsto x^{*}\) 是从 E 到其对偶上的一个双射的、半线性的映射（关于自同构 \(\xi \mapsto \overline{\xi}\) ），即从 E 到 \(E'\) 上，并且是从赋范空间 E 到赋范空间 \(E'\) 上的一个等距。

映射 \(x \mapsto x^*\) 由 (2)（V, p.~1）是半线性的，且由 Cauchy-Schwarz 不等式，我们有 \(\| x^* \| = \sup_{\| y \| \leqslant 1} |\langle x | y \rangle| = \| x \|\) ，故 \(x \mapsto x^*\) 是从 E 到 E' 中的一个等距，特别地是单射。为完成证明，需要证明对 E' 中所有 \(x' \neq 0\) ，存在 \(x \in E\) 使 \(x' = x^*\) 。而超平面 \(H = Ker x'\) 在 E 中是闭的；它的正交是一条直线 D。设 \(b\) 为 D 的一个非零元素；线性型 \(b^*\) 的核等于 H，因此存在标量 \(\lambda \neq 0\) 使 \(x' = \lambda . b^* = (\overline{\lambda}. b)^*\) 。证毕。

映射 \(x \mapsto x^{*}\) （从 E 到其对偶 \(E'\) 上的）称为典范映射。从 \(E'\) 到 E 上的逆映射也称为典范的，记作 \(x' \mapsto x'^{*}\) 。我们有

\[
\langle x | y \rangle = \langle y, x ^ {*} \rangle , \quad \langle x, x ^ {\prime} \rangle = \langle x ^ {\prime *} | x \rangle\tag{17}
\]

对 E 中的 \(x, y\) 与 E' 中的 \(x'\) 。还有 \((x^{*})^{*} = x\) 对 \(x \in \mathbf{E}\) 成立。

当 K 是 R 时，映射 \(x \mapsto x^{*}\) 是线性的。我们将用这个映射把 E 的内积转移到 \(E'\) 上。当 \(K = C\) 时，可以把映射 \(x \mapsto x^{*}\) 看作从向量空间 \(\overline{E}\) （E 的共轭）到 \(E'\) 上的同构（V, p.~6）。我们将用这个映射把 \(\overline{E}\) 的内积转移到 \(E'\) 上。

在所考虑的两种情形中，\(E'\) 是希尔伯特空间，且我们有公式

\[
\langle x ^ {*} | y ^ {*} \rangle = \overline {{\langle x | y \rangle}}, \quad \langle x ^ {\prime} | x ^ {\prime} \rangle = \| x ^ {\prime} \| ^ {2}
\]

对 \(x, y\) （\(\mathbf{E}\) 中的）与 \(x'\) （\(\mathbf{E}'\) 中的）。

说向量 \(x \in E\) 与向量 \(y \in E\) 正交，等价于说线性型 \(x^{*} \in E'\) 在 II, p.~41 所定义的意义下与 y 正交（这证明在两种情形下都可用«正交»一词）。若 M 是 E 的闭向量子空间，则与 M 正交的子空间 \(M^{\circ}\) （\(E'\) 中的，II, p.~44）是 E 中 M 的正交（定义于 V, p.~13）在 \(x \mapsto x^{*}\) 下的像（这证明在两种情形下都可用记号 \(M^{\circ}\)）。

推论 1. --- 为使希尔伯特空间 E 的点族 \((x_{i})_{i\in I}\) 是完全的，必要且充分的是：关系式 \(\langle x_{i}|y\rangle = 0\) （对 \(y \in E\) 且对所有指标 \(i \in I\) ）蕴含 y = 0。

事实上，这就是说 0 是 \(\mathbf{E}'\) 中与所有 \(x_{i}\) 正交的唯一的向量（II, p.~43 与 IV, p.~1）。

推论 2. --- 设 E 与 F 为两个希尔伯特空间。对 \(u \in \mathcal{L}(\mathrm{E};\mathrm{F})\) 、\(x \in \mathrm{E}\) 与 \(y \in \mathrm{F}\) ，令

\[
\Phi_ {u} (y, x) = \langle y | u (x) \rangle .\tag{18}
\]

映射 \(u \mapsto \Phi_u\) 是从 Banach 空间 \(\mathcal{L}(\mathrm{E};\mathrm{F})\) 到所有连续半双线性型 \(^1\) （\(\mathrm{F} \times \mathrm{E}\) 上的）组成的空间（赋以范数

\[
\| f\| = \sup_{\substack{x\in \mathbf{E},  y\in \mathbf{F}\\ \| x\| \leqslant 1, \| y\| \leqslant 1}}\left|f(y,x)\right|.\tag{19}
\]

显然 \(\Phi_u\) 对所有 \(u \in \mathcal{L}(\mathrm{E};\mathrm{F})\) 是半双线性的且连续的。反之，设 \(f\) 为 \(\mathrm{F} \times \mathrm{E}\) 上的连续半双线性型。对每个 \(x \in \mathrm{E}\) ，映射 \(y \mapsto \overline{f(y,x)}\) 是希尔伯特空间 \(\mathrm{F}\) 上的连续线性型。由定理 3，对每个 \(x \in \mathrm{E}\) ，存在唯一的元素 \(u(x)\) （在 \(\mathrm{F}\) 中）使 \(\overline{f(y,x)} = \langle u(x)|y\rangle\) 对所有 \(y \in \mathrm{F}\) 成立。映射 \(u: x \mapsto u(x)\) （从 \(\mathrm{E}\) 到 \(\mathrm{F}\) 中）是线性的，且我们有

\[
\begin{array}{l} \| f \| = \sup _ {\| x \| \leqslant 1} \sup _ {\| y \| \leqslant 1} | f (y, x) | = \sup _ {\| x \| \leqslant 1} \sup _ {\| y \| \leqslant 1} | \langle y | u (x) \rangle | \\ = \sup _ {\| x \| \leqslant 1} \| u (x) \|; \end{array}
\]

于是 \(u\) 属于 \(\mathcal{L}(\mathrm{E};\mathrm{F})\) ，\(f = \Phi_u\) 且 \(\| u\| = \| f\|\) 。这就证明了推论 2。

\(^{1}\) 回顾（A, IX, § 1, No.~5），F × E 上的（左）半双线性型 f 是从 F × E 到 K 中的映射，且满足 V, p.~1 的关系式 (1) 与 (2)。

从 E 到其双对偶 \(E''\) 的典范映射（IV, p.~14）把 E 映到 \(E''\) 上，换句话说（IV, p.~16），E 是自反 Banach 空间。事实上，若 E 是实（相应地，复）希尔伯特空间，则典范映射 \(\phi\) （从 \(E'\) 到 E 上的）是从赋范空间 \(E'\) 到 E（相应地，到 E 的共轭空间 \(\overline{E}\) ）上的一个同构；对 E（相应地，\(\overline{E}\) ）应用定理 3，我们看到赋范空间 \(E'\) 上的每个连续线性型都具有形式 \(x' \mapsto \langle \phi(x') | x \rangle = \langle x, x' \rangle\) ，其中 \(x \in E\) ，由此得出我们的断言。

作为推论（IV, p.~17, 命题 6）：

定理 4. --- 在希尔伯特空间 E 中，单位球是弱紧的。

命题 10. --- 若在希尔伯特空间 E 中，滤子 F 弱收敛于 \(x_{0}\) ，且又有 \(\lim_{\mathfrak{F}} \|x\| = \|x_{0}\|\) ，则 F 对 E 的初始拓扑收敛于 \(x_{0}\) 。事实上，\(\|x - x_{0}\|^{2} = \|x\|^{2} - 2R \langle x|x_{0} \rangle + \|x_{0}\|^{2}\) 。由于按假设 \(\langle x|x_{0}\rangle\) 关于 F 趋于 \(\|x_{0}\|^{2}\) ，且 \(\|x\|\) 关于 F 趋于 \(\|x_{0}\|\) ，故 \(\|x - x_{0}\|\) 关于 F 趋于 0，由此得命题。

注. --- 若 E 是 Hausdorff 准希尔伯特空间，\(\hat{E}\) 是 E 的希尔伯特空间完备化，我们知道（III, p.~16）E 的对偶 \(E'\) 可与 \(\hat{E}\) 的对偶等同；于是由定理 3（V, p.~15），E 上的每个连续线性型都可以唯一地写成 \(x \mapsto \langle a|x\rangle\) ，其中 \(a \in \hat{E}\) 。

